\chapter{The Ecosystem, the Alternatives and the Criticism}
\label{ch:ecosystem}

\begin{quote}\itshape
The most useful thing you can say about a framework is where it is the wrong
choice.
\end{quote}

\noindent
Sixteen chapters of one library is enough to lose perspective. This chapter is
the correction: what else exists, what those alternatives are genuinely better
at, where no framework at all is the right answer, and what the well-founded
criticisms of LangChain are --- as distinct from the ones that were true in 2023
and have been repeated without checking ever since.

\chapterstub{
\textbf{Argument.} For one or two model calls behind a clean function, the
provider SDK is faster to write, easier to debug and less likely to break. The
framework earns its place at the point where you need durability, human
interruption, or a topology --- and not before.

\textbf{Sections.} (16.1) MCP and the adapters: loading tools from a server,
decoupling tool execution from agent logic, and how this becomes the answer to
``how do you scale to hundreds of tools''. (16.2) Deep Agents: planning,
subagents, a virtual filesystem --- and the fact that it returns a compiled
graph, so nothing is a one-way door. (16.3) The alternatives, each with what it
is actually better at: PydanticAI, the OpenAI Agents SDK, CrewAI, AutoGen.
(16.4) Microsoft Agent Framework as the bridge from a .NET background: the
concepts that map one-to-one, the ones that do not, and what each ecosystem
does better --- written from having used both. (16.5) Protocols: A2A and AG-UI.
(16.6) When not to use a framework, with the decision stated as a short list of
questions. (16.7) The credible criticism: abstraction depth against debugging
cost, documentation churn, the migration tax the 1.0 split imposed, dependency
surface, and the gap between what a trace shows and what actually happened.
Written to be defensible, not contrarian --- and paired with what 1.0 genuinely
fixed.

\textbf{Listings.} An MCP server's tools loaded and given to an agent; the same
small task in LangGraph, in the provider SDK alone, and in one alternative
framework, with line counts and a note on what each makes hard.

\textbf{Diagrams.} \code{eco-decision-tree} --- framework or not, and which,
as a flowchart. \code{eco-mcp-topology} --- tools behind MCP servers rather than
in-process. \code{eco-maf-langgraph} --- the concept mapping between Microsoft
Agent Framework and LangGraph, side by side.

\textbf{Source topics covered.} Part~VIII 8.1--8.7. 8.7 is a source gap and is
the section that most needs the author's own judgement rather than research.
}
